% !PRINTED 254
\weberpage{254}

Wenn nun $x_1$ eine reelle Wurzel von (1) ist, so ist $y_1^2$ ein
positives Quadrat in der Summe (4). Sind aber $x_1$ und $x_2$ conjugirt
 imaginär, so zerfallen auch $y_1$ und $y_2$ in zwei conjugirt
imaginäre Bestandtheile
\[
y_1 = u + vi, \quad y_2 = u - vi,
\]
wo $u$ und $v$ lineare reelle Functionen von den $t$ sind. Demnach
wird
\[
y_1^2 + y_2^2 = 2u^2 - 2v^2.
\]

Es ist also durch (4) die durch (5) definirte Function
$\Phi(t)$ in eine Summe von Quadraten zerlegt, deren Anzahl
% The source visibly letterspaces only this occurrence of "verschiedenen Wurzeln".
gleich der Zahl der \gesperrt{verschiedenen Wurzeln} von $f(x) = 0$
ist, und unter denen so viele negative sind, als unter diesen
von einander verschiedenen Wurzeln Paare imaginärer
Wurzeln vorkommen.

Diese Anzahlen bleiben aber nach dem Trägheitsgesetz der
quadratischen Formen (\S.\,58) bei jeder anderen linearen reellen
Transformation der quadratischen Form in eine Summe von
Quadraten dieselben. Nun ist $S(y)$ eine reelle lineare Function
der $t$, und $B$ enthält die Variable $t_{n-1}$ nicht mehr; hiernach ergiebt
 sich aus der Formel (5) der Satz:

\gesperrt{I. Wenn die Bezoutiante $B$ durch reelle lineare
Transformation in eine Summe von $\pi$ positiven
und $\nu$ negativen Quadraten, die nicht auf eine
kleinere Zahl reducirt werden können, zerlegt
ist, so ist $\pi + \nu + 1$ die Zahl der verschiedenen
Wurzeln von $f(x) = 0$; und $\nu$ ist die Anzahl der
darunter enthaltenen Paare conjugirt imaginärer
Wurzeln, $\pi - \nu + 1$ die der reellen.}

Ist die Determinante von $B$, also auch die Discriminante
von $f(x)$ (\S.\,73) von Null verschieden, so ist $\pi + \nu = n - 1$
und $\nu$ ist kleiner oder höchstens gleich $\frac{1}{2}n$.

Die Bezoutiante ist also eine quadratische Form von einer
\gesperrt{besonderen Natur}, die sich eben darin ausspricht, dass unter
den Quadraten, in die sie sich zerlegen lässt, höchstens $\frac{1}{2}n$ negative
 vorkommen können.
